La detecció del sexisme és una preocupació creixent en el processament del llenguatge natural (PLN), però sabem poc sobre els patrons per identificar-lo. Comprendre els patrons pot aportar informació sobre els processos cognitius i millorar la interpretabilitat dels models.

En aquest treball comparem la lectura humana amb la prominència atribuïda pels models. Persones i models se centren en regions coincidents però diferents, i les persones són més sensibles als indicis pragmàtics del sexisme implícit.

Dissenyem un experiment d'eye-tracking amb 10 participants que llegeixen texts en castellà d'un corpus de sexisme. Els participants emeten judicis binaris, valoren la confiança i participen en entrevistes de record estimulat. Paral·lelament, extraiem explicacions de models de llenguatge codificadors.

Correlacionem les durades de fixació amb les atribucions dels models i comparem les fonts amb les anotacions de referència. En fonamentar l'estudi de la interpretabilitat en dades cognitives, aquest treball connecta la psicolingüística i el PLN.

\vspace{0.5cm}
\vspace{0.5cm}

\noindent{{\bfseries \large Paraules clau:} {Detecció de sexisme, eye-tracking, PLN, interpretabilitat}} \vspace{5mm}

\noindent{\bfseries \large Abstract: }

\noindent{
Sexism detection is a growing concern in NLP, yet little is known about how humans and language models allocate attention when identifying it. Understanding these patterns can illuminate cognitive processes, improve interpretability, and inform annotation guidelines.

In this work, we compare human reading behavior with model attention mechanisms. We hypothesize that humans and models attend to partially overlapping but distinct textual regions, with humans showing greater sensitivity to pragmatic cues, particularly in implicit sexism.

We design an eye-tracking experiment with 10 gender-balanced participants reading Spanish texts from MuSeD, a multimodal sexism dataset. Participants provide binary judgments and confidence ratings, followed by stimulated recall interviews. Computationally, we extract attention from encoder language models.

We correlate fixation durations with attention weights, evaluating both against gold-standard sexist spans. By grounding interpretability in cognitive data, this work bridges psycholinguistics and NLP, with implications for bias detection.
} \vspace{0.5cm}

\noindent{{\bfseries \large Keywords:} {Sexism detection, eye-tracking, NLP, interpretability}} \vspace{5mm}
